\documentclass[10pt]{article}
\usepackage[T1]{fontenc}
\usepackage[utf8]{inputenc}
\usepackage[french]{babel}
\usepackage{amsmath,amssymb,mathrsfs,mathpazo}
\usepackage{geometry}
\geometry{a4paper,margin=2cm}
\usepackage{graphicx,array,tabularx,tabularray,colortbl,multirow,multicol,booktabs,enumitem,pifont,ulem,fancybox,diagbox,cancel}
\usepackage{pstricks,pstricks-add,pst-plot,pst-tree,pst-node,pst-text,pst-func,pst-math,pst-bspline,pst-3dplot,pst-eucl,pst-solides3d}
\let\pscalculate\undefined
\usepackage{pst-bezier}
\usepackage{tikz,pgfplots,tkz-tab}
\usetikzlibrary{arrows,arrows.meta,patterns,calc,positioning,trees,decorations.markings,shapes.geometric,shapes.misc,shapes.arrows,intersections,decorations.pathmorphing,decorations.pathreplacing,backgrounds,shadows,babel}
\pgfplotsset{compat=1.18}
\usepgfplotslibrary{fillbetween}
\usepackage[np]{numprint}
\usepackage{listings,makecell,pas-tableur,twoopt}
\let\barre\undefined
\usepackage[pstricks]{tablvar}
\definecolor{darkgreen}{RGB}{0,100,0}
\providecommand{\red}{\color{red}}
\providecommand{\blue}{\color{blue}}
\providecommand{\green}{\color{green}}
\providecommand{\black}{\color{black}}
\providecommand{\white}{\color{white}}
\providecommand{\gray}{\color{gray}}
\providecommand{\cyan}{\color{cyan}}
\providecommand{\magenta}{\color{magenta}}
\providecommand{\yellow}{\color{yellow}}
\providecommand{\violet}{\color{violet}}
\providecommand{\orange}{\color{orange}}
\providecommand{\euro}{\texteuro}
\providecommand{\EUR}{\texteuro}
\providecommand{\ieme}{\textsuperscript{e}}
\providecommand{\degres}{\ensuremath{^\circ}}
\providecommand{\degree}{\ensuremath{^\circ}}
\providecommand{\vv}[1]{\overrightarrow{#1}}
\providecommand{\Coord}[1]{#1}
\providecommand{\decosix}{}
\providecommand{\starredbullet}{\ensuremath{\bullet}}
\pagestyle{empty}
\setlength{\parindent}{0pt}
\usepackage{adjustbox}\DeclareUnicodeCharacter{2236}{:}\DeclareUnicodeCharacter{2190}{\ensuremath{\leftarrow}}
\providecommand{\R}{\mathbb{R}}
\providecommand{\N}{\mathbb{N}}
\providecommand{\Z}{\mathbb{Z}}
\providecommand{\Q}{\mathbb{Q}}
\providecommand{\D}{\mathbb{D}}
\providecommand{\C}{\mathbb{C}}
\providecommand{\e}{{\mathrm{e}}}
\providecommand{\vect}[1]{\overrightarrow{#1}}
\providecommand{\vectt}[1]{\overrightarrow{\mathrm{#1}}}
\providecommand{\vvt}[1]{\overrightarrow{\mathrm{#1}}}
\providecommand{\np}[2][]{{#2~#1}}
\providecommand{\vv}[1]{{\overrightarrow{#1}}}
\renewcommand{\d}{{\mathrm{d}}}
\providecommand{\degres}{{^\circ}}
\providecommand{\degre}{{^\circ}}
\providecommand{\blue}{\color{blue}}
\providecommand{\red}{\color{red}}
\providecommand{\enskip}{\quad}
\providecommand{\llbracket}{\lbrack\!\lbrack}
\providecommand{\rrbracket}{\rbrack\!\rbrack}
\providecommand{\euro}{{\text{€}}}
\providecommand{\no}{{\text{n°}}}
\providecommand{\interval}[2]{{[#1;#2]}}
\providecommand{\intervFO}[2]{{[#1;#2[}}
\providecommand{\dotfill}{\ldots}

\providecommand{\ds}{\displaystyle}
\providecommand{\textding}[1]{#1}
\providecommand{\Oij}{\(\mathrm{O};\vec{\imath},\vec{\jmath}\)}
\providecommand{\Oijk}{\(\mathrm{O};\vec{\imath},\vec{\jmath},\vec{k}\)}
\providecommand{\Ouv}{\(\mathrm{O};\vec{u},\vec{v}\)}
\providecommand{\intervFF}[2]{\ensuremath{[#1;#2]}}
\renewcommand{\euro}{\ifmmode\text{\texteuro}\else\texteuro\fi}
\DeclareUnicodeCharacter{00B0}{\ifmmode^{\circ}\else\textdegree\fi}
\DeclareUnicodeCharacter{20AC}{\ifmmode\text{\texteuro}\else\texteuro\fi}
% Énoncés : Éducation nationale. Transcription et corrigés : APMEP, Denis Vergès et contributeurs. Découpage et classement : CoopMaths, Sébastien Lozano, Éric Elter et contributeurs. Adaptation de la présentation : Maths Explorer.
% CC BY-SA : https://coopmaths.fr/www/about/
% https://coopmaths.fr/alea/static/e3c/2020/tex/e3c_2020_00_sujet3_3_cor.tex
\begin{document}
\noindent\textit{APMEP / CoopMaths. CC BY-SA. Présentation adaptée par Maths Explorer.}\par\medskip
\section*{Exercice 3 (5 points)}
	

	\subsection*{1. Justifier que la superficie de l’enclos, en m$^2$, est donnée en fonction de $x$ par $g(x) = 4xe^{-0,5x}$ pour $x$ dans l’intervalle $[0; 5]$.}
	
	Pour $x \in [0; 5]$, on a $OA = x$ et $OC = f(x) = 4e^{-0,5x}$.\\ On a donc
	\[
	A(OABC) = x \times f(x) = 4xe^{-0,5x}
	\]
	
	\subsection*{2. La fonction $g$ est dérivable sur $[0; 5]$. Montrer que, pour tout réel $x$ de l’intervalle $[0; 5]$, on a $g'(x) = (4 - 2x)e^{-0,5x}$.}
	
	Pour $x \in [0; 5]$, on a
	\[
	g'(x) = 4e^{-0,5x} + 4xe^{-0,5x} \times (-0,5) = 4e^{-0,5x} - 2xe^{-0,5x} = (4 - 2x)e^{-0,5x}
	\]
	
	\subsection*{3. En déduire le tableau de variations de la fonction $g$ sur $[0; 5]$.}
	
	On sait que quel que soit le réel $a$, $e^a > 0$ ; le signe de $g'(x)$ est donc celui de $4 - 2x$.	
 $4 - 2x > 0 \iff 4 > 2x \iff x < 2 $\\
	Sur l’intervalle $[0; 2]$, la dérivée est positive, donc la fonction $g$ est croissante de $g(0) = 0$ à $g(2) = 8e^{-1} \approx 2,943$.
 $g(2) = 8e^{-1}$ et $g(5) = 20e^{-2,5} \approx 1,642$.

	
	Le tableau de variation est donc :
\begin{center}
		\begin{tikzpicture}[double distance=2pt]
		\tkzTabInit{$x$/1,$f'(x)$/1,$f(x)$/2}{$0$,$2$, $5$}
		\tkzTabLine{,+,z, -}
		\tkzTabVar{-/$0$,+/ $8\e^{-1}$ /, -/$20\e^{-2,5}$}
	\end{tikzpicture}
\end{center}

	
	\subsection*{4. Où doit-on placer le point $A$ sur $[OD]$ pour obtenir une superficie d’enclos maximale ? Donner la superficie maximale possible en arrondissant le résultat au dm$^2$.}
	
	D’après la question précédente l’enclos aura une surface maximale pour $x = 2$ et on a vu que $g(2) \approx 2,943$, soit 2,94 m$^2$ au dm$^2$ près.

\end{document}
